\sect{Large Language Models in Robotics}{llminrobotics}

LLMs have recently emerged as powerful enablers of language-driven reasoning and adaptive task execution in robotics \cite{zeng2023large,wang2025large,yoshikawa2023large}. By transforming unstructured NL instructions into structured, executable representations, LLMs bridge the gap between human intent and robotic control. Their transformer-based architectures capture long-range dependencies, contextual relations, and compositional semantics, allowing them to infer task hierarchies and parameter constraints directly from user input \cite{zeng2023large,wang2025large}. This capability enables intuitive, instruction-based control of robots without relying on predefined command grammars or vendor-specific programming interfaces.

Several architectures illustrate the integration of LLMs into robotic planning and execution pipelines. PaLM-SayCan couples a language model with an affordance-based feasibility module to decompose high-level user instructions into ranked, executable sub-tasks grounded in the robot’s available skills \cite{zeng2023large}. Similarly, RT-2 extends this paradigm by translating visual and textual inputs into generalisable action tokens that can be executed across diverse robotic platforms \cite{wang2025large}. These systems demonstrate how LLMs can act as cognitive planners that interpret goals, select tools, and generate structured task sequences, effectively linking NL reasoning to physical action.

The integration of LLMs into robotic systems offers several functional advantages. NL interfaces simplify human-robot communication by allowing non-expert users to define complex tasks through everyday expressions. Context-aware reasoning enables LLMs to adapt instructions to environmental or procedural constraints, while knowledge grounding aligns linguistic representations with the robot’s operational capabilities and world model \cite{zeng2023large,yoshikawa2023large}. Together, these properties make LLMs valuable intermediaries between symbolic reasoning and continuous control, enhancing flexibility and reducing the need for manual task scripting.

Despite their potential, several technical challenges remain. Reliable grounding of language to perception and actuation is limited by data availability and domain specificity. Model hallucination and overgeneralisation can lead to unsafe or infeasible plans, necessitating verification layers and guardrails within the control loop \cite{yoshikawa2023large}. The high computational demand of large models constrains real-time deployment, motivating research into smaller, locally hosted architectures for embedded robotic systems \cite{wang2025large}. Addressing these limitations is essential for achieving trustworthy, domain-adapted language-driven autonomy.

These developments establish LLMs as high-level reasoning engines capable of translating human objectives into executable robotic actions. Building on this foundation, the next stage of research focuses on LLM implementations that interact with physical systems through structured function calls, forming the conceptual basis for agentic control in autonomous laboratory environments.
